\chapter{Conclusioni}
Lo scopo di questo lavoro di tesi è quello di permettere ad un utente di controllare remotamente il proprio dispositivo attraverso l'utilizzo del linguaggio naturale. 
Per raggiungere tale obiettivo, è stato utilizzato l'assistente virtuale Google Assistant per tradurre le frasi in concetti comprensibili da una applicazione sul dispositivo che si vuole controllare. 
Per quanto riguarda l'attuale stato dell'applicazione, risulta evidente che è basato su troppi software esterni, ma è possibile rendere il tutto una esperienza più naturale e coesa implementando alcuni dei miglioramenti descritti nel capitolo precedente. 
A causa della natura sperimentale di quello che si voleva realizzare, si è finiti con il prioritizzare la comodità del programmatore a quella dell'utente finale, cosa non neccessariamente sbagliata, ma nemmeno la più corretta. 
Fortunatamente, l'informatica è in continua evoluzione, e questo potenziale rende possibile migliorare qualsiasi aspetto necessiti di cambiamenti.